\documentclass[12pt,a4paper,titlepage]{article}
\usepackage[utf8]{inputenc}
\usepackage[spanish]{babel}
\usepackage{graphicx}
\usepackage[hidelinks]{hyperref}
\usepackage{geometry}
\geometry{a4paper, margin=2.5cm}
\usepackage{titlesec}
\newcommand{\sectionbreak}{\clearpage}

\title{Manual de Operación y Automatización del Observatorio\\
\large Basado en RTS2 (Remote Telescope System 2)}
\author{Universidad Politécnica de Madrid \\ Escuela Técnica Superior de Ingenieros Informáticos (ETSIINF)}
\date{\today}

\begin{document}

\maketitle
\tableofcontents
\newpage

\section{Introducción}
El presente documento constituye el Manual Oficial de Operación y Automatización del observatorio astronómico de la Escuela Técnica Superior de Ingenieros Informáticos (ETSIINF) de la Universidad Politécnica de Madrid (UPM).

Este proyecto nace de la iniciativa de un grupo de estudiantes con el objetivo de reparar, actualizar y automatizar las instalaciones del observatorio, el cual cuenta como instrumento principal con un telescopio \textbf{Meade LX200} y una cúpula motorizada.

Para llevar a cabo esta automatización, se ha adoptado el software \textbf{RTS2} (Remote Telescope System 2), un entorno de código abierto ampliamente utilizado a nivel mundial para el control robótico de observatorios astronómicos. RTS2 permite no solo coordinar el movimiento del telescopio y la apertura de la cúpula, sino también programar secuencias de observación desatendidas, gestionar cámaras y registrar condiciones meteorológicas.

El objetivo principal de esta guía es proporcionar a los futuros alumnos e investigadores de la facultad un recurso claro, técnico y estructurado. Aquí se detalla desde la configuración inicial en sistemas Linux y la compilación del código fuente, hasta los comandos diarios de operación y la resolución de problemas (troubleshooting).

\section{Arquitectura del Sistema RTS2}
RTS2 (Remote Telescope System 2) está diseñado como un sistema distribuido y modular basado en \textit{demonios} (procesos que se ejecutan en segundo plano en Linux). Esto permite aislar los fallos: si la cámara falla, el demonio de la cámara se reinicia sin afectar al seguimiento del telescopio.

Todos los demonios se comunican entre sí enviando mensajes a través de un puerto de red local (habitualmente el puerto 8889).

\subsection{Componentes y Demonios Principales}
En nuestro observatorio, la arquitectura se compondrá principalmente de los siguientes elementos:

\begin{itemize}
    \item \textbf{centrald}: Es el cerebro del sistema y el enrutador de mensajes. Conoce el estado global y permite que los demás demonios hablen entre ellos. Debe ejecutarse siempre el primero.
    \item \textbf{teld} (Telescope Daemon): Controla el apuntado, el seguimiento y el movimiento de la montura. Para nuestro telescopio, utilizaremos el módulo específico de \textbf{Meade LX200}.
    \item \textbf{dome} (Dome Daemon): Controla el movimiento de la cúpula, asegurándose de que la apertura (la rendija) esté siempre alineada con donde apunta el telescopio.
    \item \textbf{camd} (Camera Daemon): Se encarga de controlar la cámara principal, gestionar los tiempos de exposición y la refrigeración.
    \item \textbf{execd} (Execution Daemon): Es el encargado de leer los planes de observación (qué observar, cuándo y cómo) y enviar las órdenes en secuencia al resto del sistema.
    \item \textbf{Base de Datos (PostgreSQL)}: RTS2 guarda todos los objetivos astronómicos, registros de observación y logs en una base de datos central para poder automatizar el observatorio.
\end{itemize}

\section{Configuración del Entorno (Linux)}
Los archivos de configuración principales de RTS2 se alojan en el directorio \texttt{/etc/rts2/}. 

\subsection{Archivo \texttt{rts2.ini}}
(En esta sección se documentará el contenido del archivo principal extraído del ordenador del laboratorio).

\subsection{Configuración del Telescopio (Meade LX200)}
El Meade LX200 se conecta típicamente a través de un puerto serie o adaptador USB-Serie (ej: \texttt{/dev/ttyUSB0}).
(Aquí añadiremos las líneas de configuración exactas una vez tengamos los archivos).

\section{Arranque y Operación}
\subsection{Secuencia de Inicio}
Para iniciar el sistema RTS2 de forma manual en el terminal, se debe respetar el estricto orden de dependencia de los demonios. El orden estándar es:
\begin{verbatim}
$ rts2-centrald
$ rts2-teld-lx200 -d /dev/ttyUSB0
$ rts2-dome
\end{verbatim}
\textit{Nota: En un entorno de producción final, estos servicios se configurarán para que inicien automáticamente en el arranque del sistema usando \texttt{systemd} o scripts init.}

\subsection{Interfaz de Usuario (\texttt{rts2-mon})}
Para interactuar con el observatorio, ver el estado de los componentes y mandar comandos de forma manual, RTS2 incluye una potente consola interactiva llamada \texttt{rts2-mon}. 

Al ejecutar \texttt{rts2-mon} en la terminal, se despliega una interfaz basada en texto (ncurses) donde cada dispositivo (Telescopio, Cúpula, Cámara) tiene su propia pestaña. Usando las flechas del teclado, los usuarios pueden navegar y enviar comandos en tiempo real, como por ejemplo introducir las coordenadas de Ascensión Recta (RA) y Declinación (DEC) y ordenarle al Meade LX200 que comience a moverse hacia el objetivo (\textit{slewing}).

\section{Solución de Problemas Frecuentes (Troubleshooting)}

Esta sección irá creciendo a medida que nos encontremos con bugs y casuísticas reales del montaje, pero los problemas iniciales más comunes a comprobar son:

\begin{itemize}
    \item \textbf{El telescopio no responde a los comandos de movimiento:}
    \begin{itemize}
        \item Verificar que el cable USB/Serie está bien conectado al PC y detectar qué puerto se le ha asignado ejecutando \texttt{dmesg | grep tty}.
        \item Verificar los permisos de lectura/escritura del puerto serie (ej: \texttt{sudo chmod a+rw /dev/ttyUSB0}).
        \item Comprobar si el demonio \texttt{teld} sigue ejecutándose o si ha crasheado ejecutando \texttt{ps aux | grep rts2-teld}.
        \item Revisar los logs específicos del telescopio guardados en \texttt{/var/log/rts2/} o en \texttt{/var/log/syslog}.
    \end{itemize}
    
    \item \textbf{Error de conexión al iniciar los demonios (no encuentran \texttt{centrald}):}
    \begin{itemize}
        \item Esto ocurre si \texttt{centrald} no está corriendo o si se cerró por algún error. Hay que reiniciar todos los servicios asegurándose de que \texttt{centrald} se levante primero.
    \end{itemize}
    
    \item \textbf{Fallo general de RTS2 al registrar datos de observación:}
    \begin{itemize}
        \item RTS2 requiere que el motor de base de datos PostgreSQL esté activo. Se puede comprobar su estado en Linux ejecutando \texttt{systemctl status postgresql}.
    \end{itemize}
\end{itemize}

\end{document}
